% Generated draft: requires mathematical and source review.
\documentclass[11pt]{article}
\usepackage[margin=25mm]{geometry}
\usepackage{fontspec}
\setmainfont{Latin Modern Roman}
\usepackage{amsmath,amssymb,mathtools,unicode-math}
\setmathfont{Latin Modern Math}
\usepackage{enumitem,needspace,xurl,hyperref}
\setlength{\parindent}{0pt}
\setlength{\parskip}{5pt}
\setlength{\emergencystretch}{3em}
\newcommand{\N}{\mathbb{N}}
\newcommand{\Z}{\mathbb{Z}}
\newcommand{\Q}{\mathbb{Q}}
\newcommand{\R}{\mathbb{R}}
\newcommand{\C}{\mathbb{C}}
\newcommand{\F}{\mathbb{F}}
\newcommand{\A}{\mathbb{A}}
\newcommand{\PP}{\mathbb{P}}
\newcommand{\E}{\mathbb{E}}
\newcommand{\eps}{\varepsilon}
\newcommand{\dd}{\,\mathrm{d}}
\begin{document}
\section*{UnsolvedMath ID: 11100063}
\textbf{Source problem code:} catalogue entry
\subsection{Definitions and mathematical statement}
Let $MU$ denote the complex bordism spectrum, equivalently the Thom spectrum of the universal stable complex bundle. Let $\mathcal M_{\mathrm{fg}}$ denote the moduli stack of one-dimensional commutative formal groups, customarily over $\operatorname{Spec}\Z$: it records formal groups and their isomorphisms, rather than only formal group laws with a chosen coordinate. Base and stack-topology conventions are not stated in the excerpt.

\textbf{Miscellaneous problem 1 (Hovey).} Construct $MU$ from $\mathcal M_{\mathrm{fg}}$. The goal is a topological or homotopical construction from the stack that recovers the spectrum, not merely the algebraic description of complex-bordism coefficients or the familiar cooperation Hopf algebroid. In particular, a presentation that already uses $MU$ as input does not by itself give the requested construction. The source does not specify the allowable construction, functoriality, equivalence notion, or whether ring-spectrum and complex-orientation structure must also be recovered.

Hovey suggests treating the stack as a space in Voevodsky's homotopy-theoretic setting and asks whether it might be an infinite loop space there. This is a tentative proposed route, not an assertion that the stack already has compatible deloopings. A precise version needs an ambient motivic category, base, descent topology, localization, meaning of delooping, and a realization/comparison to topological spectra. Recovering $MU$ rather than only its formal-group algebra is the intended endpoint.

The source calls this an old problem heard in Ravenel's ``green book.'' The summary invokes a theorem of Lurie explaining constraints from the formal-group stack but provides no precise theorem or construction matching this request. Neither that report nor the speculative motivic route is asserted here as a completed recovery of $MU$.
\subsection{Short English statement}
Recover the complex bordism spectrum from the moduli stack of formal groups. The source proposes motivic infinite-loop-space ideas, but leaves the construction and realization procedure unspecified.
\subsection{Sources}
\begin{itemize}
\item Local UnsolvedMath registry, permanent ID 11100063, source code catalogue entry; attributed to Mark Hovey, no proposal year supplied.
\item \emph{Hovey's Algebraic Topology Problem List}, Miscellaneous problems, item 1; surviving mirror \url{https://www-users.cse.umn.edu/~tlawson/hovey/}. The registry records HTML-ordered-list extraction and access on 2026-07-29; the mirror has 65 top-level items on nine category pages.
\item The source names Ravenel's ``green book'' and Voevodsky's category without exact bibliographic data. The summary names a theorem of Lurie without a precise citation or statement; no title or theorem has been guessed.
\end{itemize}
\end{document}
